\section{Learning where movement changes}
\label{sec:bend-feedback}
A distance score gives the same answer to different sequences of motion. To see the problem directly, we permuted the one-second speeds within each of 22 SUMO bus references. Total travel and the empirical speed distribution in the 127 complete seconds were identical before and after permutation. Yet the average longest continuous steady-moving run fell from 29.6 seconds to 4.4. We mark a second as steady when speed is at least 0.3 m/s and the absolute change from the preceding second is below 0.5 m/s. This operational definition also allows slow steady movement.

The previous eight-class generator provided a second diagnostic. On G14, the reference buses travelled 2.40 m/s slower in curved road sections than in straight ones. The generated buses slowed by only 0.39 m/s. For each trajectory we sampled arc centres every five metres, comparing headings over the ten metres immediately before and after each centre. A heading change of at least ten degrees marked a curved section; at most three degrees marked a near-straight section. Local speed was measured over the central five metres. Each trajectory contributed its own curved-minus-straight speed difference. The reference and generator can traverse different roads. The comparison measures the association between movement and road geometry; it cannot attribute braking to a particular bend, traffic light or leader.

\paragraph{Road-conditioned diffusion.}
We introduced an explicit spatial speed effect. A training-only table estimates typical local speed by class, bend severity and posted lane speed, shrinking sparse cells towards the class average. It defines a positive factor $w_c(s)$, bounded between 0.5 and 2. Along a supplied path, physical distance $s$ corresponds to a transformed coordinate $u(s)=\int_0^s w_c(r)\,dr$. The diffusion model learns increments in this coordinate. At sampling, inverting the cumulative map gives physical progress and the usual path decoder returns timed positions. Larger values of $w_c$ yield less physical progress per unit of generated increment. The model receives only planned geometry at inference. Lane limits are read from training annotations; at inference they are approximated from the nearest compatible lane. Posted limits do not encode actual traffic signal phases.

Both comparison arms started from the same respective G17 weights and used the same training windows, class sampling, noise, optimizer and 10,000 additional updates. There were two seeds per arm and 63,717 eligible SUMO training windows. The control used ordinary physical arc increments. The change therefore measures the combined effect of the calibrated spatial transformation and training in that coordinate. G14 is a separate simulator population that had already been examined during development. Both test sets remained reserved for their original purposes.

\begin{table}[htbp]\centering\small
\caption{Two motion-pattern errors on G14. For each class we compare the generated and reference means, take the absolute difference, and then average over the five classes. Lower is better. Intervals resample the eight simulation episodes jointly, conditional on the two fits.}
\begin{tabular}{lrrr}\toprule
Physical error & Control & Road model & Paired change [95\%]\\\midrule
Bend speed gap (m/s) & 1.77 & 0.42 & $-1.35$ [$-1.63$,$-0.99$]\\
Longest steady run (s) & 18.23 & 37.55 & $+19.32$ [$+18.02$,$+20.33$]\\\bottomrule
\end{tabular}
\end{table}

The road model generated a bus slowdown of 2.76 m/s on G14, compared with 0.37 m/s for the matched control. On G15 the class-macro bend-gap error likewise fell from 1.78 to 0.45 m/s, while steady-run error increased from 7.57 to 23.97 s. All 960 G14 candidates per arm passed the directed lane audit. On G15 all vehicle candidates passed this audit and all pedestrian candidates satisfied the single-sidewalk proximity check. Four of the road model's 1,400 G15 outputs exhausted their planned path. G14 travel-distance discrepancy increased from 160.7 to 174.0 m. These measurements separate a better response to bends from fidelity to the broader sequence of motion.

\paragraph{Can inherited image feedback improve this movement?}
We extended the feedback network to eight labels, keeping the original time-ordered Hilbert sequence, GASF, GADF, MTF and frozen encoding ranges. A ResNet18 inherited the earlier self-trained scorer's backbone and fitted a new head and backbone on 8,192 class-balanced G15 training examples for five epochs. Final accuracy was 32.9\% on 1,024 internal monitor tracks. On 175 G15 development requests its accuracy was 27.4\% and macro F1 was 27.0\%. These scores describe an auxiliary recognizer built for this study.

The scorer then selected one of four diffusion candidates, with reference futures hidden until scoring was complete. On G14 it changed the choice relative to a fixed uniform slot in 180 of 240 road-model pools. All 960 candidates were encodable. The bend-gap error changed by $+0.02$ m/s (95\% interval $-0.17$ to $+0.22$), and steady-run error by $+0.56$ s (interval $-1.19$ to $+2.55$). In G15 the respective changes were $-0.03$ m/s and $+0.27$ s, with both intervals again spanning zero. The delivered model exposes all four trajectories together with the GAF choice. This makes the original representation usable and its physical effect independently measurable.
